\chapter{Контейнеры: жизненный цикл}

Контейнер живёт ровно столько, сколько живёт его главный процесс.
В этом простом правиле~--- ключ к пониманию большинства \enquote{странностей}
Docker: почему контейнер сразу завершается, почему не работают логи,
почему \cmd{docker exec} не спасает от неправильного \cmd{CMD}.
В главе разобраны состояния контейнера, команды управления, сигналы,
коды выхода и типичные сценарии отладки.

\section{Главный процесс}

При создании контейнера Docker запускает одну команду~--- она получает
PID~1 внутри контейнера. Всё остальное~--- это процессы, которые она
породила. Когда главный процесс завершается, контейнер останавливается,
даже если внутри остались дочерние процессы.

\begin{defbox}[title={Определение: контейнер (практическое)}]
Контейнер~--- это исполняемый экземпляр образа. Он состоит из
записываемого слоя файловой системы, конфигурации (сеть, тома,
переменные окружения) и главного процесса. Контейнер существует,
пока существует его главный процесс, либо пока его не удалили.
\end{defbox}

Отсюда классическая ошибка новичка: запустить \cmd{nginx} в фоне,
не удержав его на переднем плане. Контейнер тут же завершится.
Правильный запуск выглядит как \cmd{nginx -g "daemon off;"}~---
именно это и делает официальный образ \cmd{nginx}.

\begin{warnbox}
Не используйте внутри контейнера \cmd{systemd}, \cmd{cron} и другие
менеджеры служб как главный процесс. Docker сам является менеджером
служб: один контейнер~--- один процесс (или несколько тесно связанных).
Нужны две службы~--- запустите два контейнера и свяжите их сетью.
\end{warnbox}

\section{Создание и запуск}

Команда \cmd{docker run} объединяет два шага: \cmd{docker create}
(создать контейнер) и \cmd{docker start} (запустить его). Разделять
их полезно, когда контейнер нужно подготовить заранее.

\begin{code}[language=bash]
$ docker create --name web -p 8080:80 nginx:1.27-alpine
$ docker start web
$ docker ps
CONTAINER ID   IMAGE               COMMAND                  CREATED         STATUS         PORTS                  NAMES
7c9e1a2b3d4f   nginx:1.27-alpine   "/docker-entrypoint.…"   5 seconds ago   Up 3 seconds   0.0.0.0:8080->80/tcp   web
\end{code}

Основные флаги \cmd{docker run} собраны в таблице.

\begin{center}
\small
\begin{tabularx}{\textwidth}{@{}lX@{}}
\toprule
\textbf{Флаг} & \textbf{Назначение} \\
\midrule
\flag{--name ИМЯ} & задать имя вместо случайного \\
\flag{-d} & отсоединиться, работать в фоне \\
\flag{-i} & оставить stdin открытым \\
\flag{-t} & выделить псевдотерминал \\
\flag{--rm} & удалить контейнер после завершения \\
\flag{-p ХОСТ:КОНТ} & опубликовать порт \\
\flag{-e ПЕРЕМ=ЗНАЧ} & задать переменную окружения \\
\flag{--env-file ФАЙЛ} & прочитать переменные из файла \\
\flag{-v ИСТОЧНИК:ЦЕЛЬ} & подключить том или каталог \\
\flag{--network СЕТЬ} & подключить к сети \\
\flag{--restart ПОЛИТИКА} & политика перезапуска \\
\flag{--memory}, \flag{--cpus} & ограничения ресурсов \\
\flag{--user UID:GID} & запустить от указанного пользователя \\
\flag{--entrypoint КОМАНДА} & заменить \cmd{ENTRYPOINT} образа \\
\flag{--init} & добавить init-процесс для сигналов и зомби \\
\flag{--platform linux/arm64} & выбрать платформу образа \\
\bottomrule
\end{tabularx}
\end{center}

\begin{code}[language=bash]
$ docker run -d --name api \
    -p 127.0.0.1:3000:3000 \
    -e NODE_ENV=production \
    --env-file ./config/api.env \
    --memory=512m --cpus=1 \
    --restart unless-stopped \
    registry.example.com/team/api:1.4.2
\end{code}

\begin{tipbox}
Публикуйте порты на конкретном адресе (\flag{-p 127.0.0.1:3000:3000}),
если сервис не должен быть доступен из сети. Запись
\flag{-p 3000:3000} открывает порт на всех интерфейсах, включая
внешний,~--- это частая причина незапланированной публикации.
\end{tipbox}

\section{Состояния контейнера}

У контейнера несколько состояний, и переходы между ними выполняются
разными командами.

\begin{figure}[htbp]
\centering
\begin{tikzpicture}[
  font=\sffamily\scriptsize,
  state/.style={draw=accent!55!black, fill=accent!7, rounded corners=6pt,
    minimum width=2.6cm, minimum height=0.8cm},
  trans/.style={-Latex, thick, accent!65},
]
  \node[state, fill=good!10, draw=good!60!black] (running) at (4.2,2.6) {running};
  \node[state] (created) at (1.0,2.6) {created};
  \node[state, fill=warm!10, draw=warm!65!black] (paused) at (7.4,2.6) {paused};
  \node[state, fill=black!5, draw=black!40] (exited) at (4.2,0.7) {exited};
  \node[state, dashed] (removed) at (1.0,0.7) {removed};

  \draw[trans] (created) -- node[above, text=muted] {start, run} (running);
  \draw[trans] (running) -- node[above, text=muted] {pause} (paused);
  \draw[trans] (paused) -- node[below, text=muted] {unpause} (running);
  \draw[trans] ([xshift=-4pt]running.south) --
    node[left, text=muted] {stop, kill} ([xshift=-4pt]exited.north);
  \draw[trans] ([xshift=4pt]exited.north) --
    node[right, text=muted] {start} ([xshift=4pt]running.south);
  \draw[trans] (exited) -- node[above, text=muted] {rm} (removed);
  \draw[trans] (created) -- node[right, text=muted] {rm} (removed);
\end{tikzpicture}
\caption{Жизненный цикл контейнера. Политика перезапуска автоматически
переводит контейнер из \cmd{exited} обратно в \cmd{running}}
\label{fig:lifecycle}
\end{figure}

\begin{itemize}
  \item \cmd{created}~--- контейнер создан, но процесс не запущен;
  \item \cmd{running}~--- главный процесс работает;
  \item \cmd{paused}~--- процессы заморожены (cgroup freezer), память
        сохраняется; используется редко, для отладки и миграций;
  \item \cmd{exited}~--- процесс завершился; файловая система и логи
        сохраняются до удаления;
  \item \cmd{removed}~--- контейнер удалён; записываемый слой уничтожен.
\end{itemize}

Посмотреть состояние можно в \cmd{docker ps -a} или в JSON:

\begin{code}[language=bash]
$ docker inspect --format '{{.State.Status}} exit={{.State.ExitCode}}' web
running exit=0
\end{code}

\section{Остановка, сигналы и коды выхода}

Когда вы выполняете \cmd{docker stop web}, Docker посылает главному
процессу \cmd{SIGTERM}, ждёт 10~секунд и, если процесс не завершился,
посылает \cmd{SIGKILL}. Приложение должно корректно обрабатывать
\cmd{SIGTERM}: закрывать соединения, дописывать файлы, освобождать
блокировки.

\begin{code}[language=bash]
$ docker stop web                 # SIGTERM, затем SIGKILL через 10 с
$ docker stop -t 30 web           # ждать 30 секунд
$ docker kill web                 # сразу SIGKILL
$ docker kill -s SIGHUP web       # послать конкретный сигнал
\end{code}

Коды выхода помогают понять причину остановки.

\begin{center}
\small
\begin{tabularx}{\textwidth}{@{}lX@{}}
\toprule
\textbf{Код} & \textbf{Значение} \\
\midrule
\cmd{0} & успешное завершение \\
\cmd{1} & ошибка приложения (смотрите логи) \\
\cmd{126} & файл найден, но не исполняется (нет прав или неверный формат) \\
\cmd{127} & команда не найдена (опечатка в \cmd{CMD} или нет оболочки) \\
\cmd{137} & процесс убит \cmd{SIGKILL} (128+9), часто из-за лимита памяти \\
\cmd{143} & процесс завершён \cmd{SIGTERM} (128+15) \\
\bottomrule
\end{tabularx}
\end{center}

\begin{warnbox}
Код \cmd{137}~--- не всегда ошибка приложения. Если контейнер убит
из-за \cmd{OOMKilled}, увеличьте лимит памяти или исправьте утечку.
Проверить причину: \cmd{docker inspect --format '\{\{.State.OOMKilled\}\}' web}
и \cmd{docker inspect --format '\{\{.State.ExitCode\}\}' web}.
\end{warnbox}

\section{Логи и стандартные потоки}

Docker перехватывает stdout и stderr главного процесса и сохраняет их
через драйвер журналирования (по умолчанию \cmd{json-file}).
Файлы, которые приложение пишет само, в логи Docker не попадают~---
их нужно читать через \cmd{docker exec} или выносить в том.

\begin{code}[language=bash]
$ docker logs web
$ docker logs -f web              # следить в реальном времени
$ docker logs --tail 100 web
$ docker logs --since 10m web
$ docker logs -t web              # с метками времени
$ docker logs web > web.log 2>&1  # сохранить в файл
\end{code}

\begin{tipbox}
Правило двенадцати факторов: приложение должно писать логи в stdout
и stderr, а не в файлы. Тогда \cmd{docker logs}, системы сбора логов
и платформы оркестрации работают без дополнительной настройки.
\end{tipbox}

\begin{warnbox}
Без ограничения размера журналы могут занять весь диск: \cmd{json-file}
растёт неограниченно. Настройте \cmd{log-opts} в \file{daemon.json}
(раздел~3.9) или задайте \flag{--log-opt max-size=10m} для контейнера.
\end{warnbox}

\section{Работа с работающим контейнером}

\cmd{docker exec} запускает дополнительный процесс внутри уже
работающего контейнера. Это главный инструмент отладки.

\begin{code}[language=bash]
$ docker exec -it web sh
$ docker exec web env
$ docker exec -u root web whoami
$ docker exec -w /usr/share/nginx/html web ls -la
\end{code}

Не путайте \cmd{exec} и \cmd{attach}:

\begin{itemize}
  \item \cmd{docker exec} создаёт \emph{новый} процесс. Оболочка
        не влияет на главный процесс и безопасно закрывается по \cmd{exit}.
  \item \cmd{docker attach} подключается к stdio главного процесса.
        Выход по \cmd{Ctrl+C} может отправить сигнал и остановить
        контейнер; отключиться без остановки~--- \cmd{Ctrl+P}, затем
        \cmd{Ctrl+Q}.
\end{itemize}

Другие полезные команды:

\begin{code}[language=bash]
$ docker cp web:/etc/nginx/nginx.conf ./nginx.conf
$ docker diff web                  # изменённые файлы в записываемом слое
$ docker top web                   # процессы контейнера
$ docker port web                  # опубликованные порты
$ docker stats --no-stream web     # потребление ресурсов
$ docker rename web web-old
\end{code}

\begin{notebox}
В минимальных образах (alpine, distroless, scratch) может не быть
\cmd{sh}, \cmd{bash}, \cmd{curl} и даже \cmd{ls}. Это нормально и полезно
для безопасности, но усложняет отладку. Используйте отладочный вариант
образа, временный контейнер с теми же томами и сетью или \cmd{docker debug}
(если он доступен в вашей подписке).
\end{notebox}

\section{Политики перезапуска}

Политика перезапуска определяет, что делать, когда главный процесс
завершился.

\begin{center}
\small
\begin{tabularx}{\textwidth}{@{}lX@{}}
\toprule
\textbf{Политика} & \textbf{Поведение} \\
\midrule
\cmd{no} & не перезапускать (по умолчанию) \\
\cmd{on-failure[:N]} & перезапускать при ненулевом коде выхода,
не более \cmd{N} раз \\
\cmd{always} & перезапускать всегда, в том числе после перезапуска демона \\
\cmd{unless-stopped} & как \cmd{always}, но не поднимать контейнер,
остановленный вручную \\
\bottomrule
\end{tabularx}
\end{center}

\begin{code}[language=bash]
$ docker run -d --restart unless-stopped --name api api:1.4.2
$ docker update --restart on-failure:3 api
$ docker update --memory 1g --cpus 2 api
\end{code}

\begin{warnbox}
Политика \cmd{always} и ошибочный \cmd{CMD} дают бесконечный цикл
перезапусков, который маскирует проблему и нагружает диск логами.
Сначала добейтесь, чтобы контейнер работал без перезапусков, и только
потом включайте политику.
\end{warnbox}

\section{Разбор типичной ситуации}

Контейнер с приложением сразу завершается. Что делать:

\begin{enumerate}
  \item \cmd{docker ps -a}~--- посмотреть код выхода и статус.
  \item \cmd{docker logs <имя>}~--- прочитать последние строки.
  \item \cmd{docker inspect <имя>}~--- проверить \cmd{Cmd}, \cmd{Entrypoint},
        \cmd{Env}, монтирования.
  \item Запустить с другой командой, чтобы получить оболочку:
        \cmd{docker run --rm -it --entrypoint sh IMAGE}.
  \item Проверить, не завершается ли главный процесс из-за отсутствия
        терминала или из-за того, что демон ушёл в фон.
\end{enumerate}

\begin{summarybox}
\begin{itemize}
  \item Контейнер живёт, пока живёт его главный процесс с PID~1.
  \item \cmd{run} = \cmd{create} + \cmd{start}; состояния
        \cmd{created}, \cmd{running}, \cmd{paused}, \cmd{exited},
        \cmd{removed}.
  \item \cmd{stop} посылает \cmd{SIGTERM}, затем \cmd{SIGKILL};
        код 137 часто означает нехватку памяти.
  \item Логи~--- это stdout и stderr; \cmd{exec} запускает новый процесс,
        \cmd{attach} подключается к главному.
  \item Политики перезапуска помогают в production, но маскируют
        ошибки при отладке.
\end{itemize}
\end{summarybox}

\section*{Вопросы для самопроверки}
\begin{selfcheck}
  \item Почему контейнер с nginx без \cmd{daemon off} сразу завершается?
  \item Чем \cmd{docker stop} отличается от \cmd{docker kill}?
  \item Что означает код выхода 137 и как проверить, не убит ли контейнер
        из-за нехватки памяти?
  \item Чем \cmd{exec} отличается от \cmd{attach}?
  \item Зачем нужна политика \cmd{unless-stopped}?
  \item Где окажутся логи, если приложение пишет их в
        \file{/var/log/app.log}?
\end{selfcheck}
